% ============================================================
% ChatGIT: RAG-Based Enhanced GitHub Repository Analysis Tool
% IEEE Conference Format — Camera-Ready
% Authors: Anish Gupta, Prakhar Sethi, Ritwik Bhattacharya, Dr Sonia Khetarpaul
% Institution: SNU (School of Engineering and Applied Science)
% ============================================================

\documentclass[conference]{IEEEtran}
\IEEEoverridecommandlockouts

\usepackage[T1]{fontenc}
\usepackage[utf8]{inputenc}
\usepackage{cite}
\usepackage{amsmath,amssymb}
\usepackage{graphicx}
\usepackage{booktabs}
\usepackage{multirow}
\usepackage{array}
\usepackage{xspace}
\usepackage{microtype}
\usepackage[table]{xcolor}
\usepackage[colorlinks=true,linkcolor=blue!60!black,citecolor=blue!60!black,
            urlcolor=blue!60!black]{hyperref}

% TikZ and pgfplots for all diagrams
\usepackage{tikz}
\usetikzlibrary{arrows.meta, shapes.geometric, shapes.rounded,
                fit, positioning, backgrounds, calc,
                decorations.pathreplacing, shadows.blur,
                fadings}
\usepackage{pgfplots}
\pgfplotsset{compat=1.18}

% ── Colour palette ────────────────────────────────────────────
\definecolor{cBlue}   {HTML}{3B82F6}
\definecolor{cIndigo} {HTML}{6366F1}
\definecolor{cPurple} {HTML}{A855F7}
\definecolor{cGreen}  {HTML}{22C55E}
\definecolor{cAmber}  {HTML}{F59E0B}
\definecolor{cRed}    {HTML}{EF4444}
\definecolor{cCyan}   {HTML}{06B6D4}
\definecolor{cSlate}  {HTML}{64748B}
\definecolor{cBg}     {HTML}{F8FAFC}
\definecolor{cDark}   {HTML}{1E293B}

% ── TikZ node styles ──────────────────────────────────────────
\tikzset{
  % Rounded pill node with drop shadow
  pnode/.style={
    rounded corners=5pt,
    draw=#1!70!black, line width=0.7pt,
    fill=#1!15,
    text=cDark,
    font=\footnotesize\bfseries,
    align=center,
    minimum height=0.65cm,
    inner sep=5pt,
    blur shadow={shadow blur steps=4, shadow xshift=0.5pt,
                 shadow yshift=-0.5pt, shadow opacity=20}
  },
  % Novelty badge (orange accent)
  nbadge/.style={
    rounded corners=3pt,
    draw=cAmber!80, line width=0.6pt,
    fill=cAmber!20,
    font=\scriptsize\bfseries,
    text=cDark,
    align=center,
    inner sep=3pt,
    minimum height=0.52cm,
  },
  % Arrow style
  arr/.style={
    -{Stealth[length=5pt, width=4pt]},
    line width=1pt,
    color=#1
  },
  % Stage label
  stlabel/.style={
    font=\scriptsize\bfseries,
    text=#1!70!black,
    rotate=90,
  },
  % Group bracket
  grpbox/.style={
    rounded corners=4pt,
    draw=#1!50, dashed, line width=0.6pt,
    fill=#1!5,
    inner sep=5pt,
  },
}

% ── Macros ────────────────────────────────────────────────────
\newcommand{\chatgit}{\textsc{ChatGIT}\xspace}
\newcommand{\nI}{$\mathcal{N}_1$}
\newcommand{\nII}{$\mathcal{N}_2$}
\newcommand{\nIII}{$\mathcal{N}_3$}
\newcommand{\nIV}{$\mathcal{N}_4$}
\newcommand{\nV}{$\mathcal{N}_5$}
\newcommand{\nVI}{$\mathcal{N}_6$}
\newcommand{\nVII}{$\mathcal{N}_7$}
\newcommand{\nVIII}{$\mathcal{N}_8$}
\newcommand{\nIX}{$\mathcal{N}_9$}

% ──────────────────────────────────────────────────────────────
\begin{document}

\title{\chatgit: A RAG-Based Conversational GitHub Repository\\
Analysis Tool with Nine Structural Retrieval Novelties}

\author{
  \IEEEauthorblockN{Anish Gupta\quad Prakhar Sethi\quad Ritwik Bhattacharya\quad Dr.~Sonia Khetarpaul}
  \IEEEauthorblockA{
    Computer Science and Engineering,
    School of Engineering and Applied Science, SNU\\
    \texttt{\{ag801, ps385, rb213\}@snu.edu.in\quad sonia.khetarpaul@snu.edu.in}
  }
}

\maketitle

% ============================================================
\begin{abstract}
Understanding large-scale software repositories via keyword search is
slow, semantically shallow, and structurally blind.
We present \chatgit, a Retrieval-Augmented Generation (RAG) system for
multi-turn conversational question-answering over GitHub repositories.
Starting from a baseline implementation (P@5\,=\,0.48, F1\,=\,0.60),
we introduce nine complementary technical novelties---spanning AST
chunking, git-volatility weighting, hybrid PageRank + query-conditioned
graph attention, multi-turn session memory, intent-driven granularity
adaptation, bidirectional call-context augmentation, cross-encoder
reranking, post-generation line-number injection, and HITS
hub/authority scoring---to reach
\textbf{P@5\,=\,0.78, Recall@10\,=\,0.86, Accuracy\,=\,0.89,
F1\,=\,0.88} on our query suite and \textbf{MRR\,=\,0.389} on
ConvCodeBench (50 repositories, 5 languages).
All nine novelties, the benchmark, and the full open-source
system are publicly released.
\end{abstract}

\begin{IEEEkeywords}
Retrieval-Augmented Generation, Code Search, AST Chunking,
Session Memory, Call Graph, PageRank, HITS, Intent Classification,
Repository Intelligence
\end{IEEEkeywords}

% ============================================================
\section{Introduction}
\label{sec:intro}

Modern software projects routinely span hundreds of files, multiple
programming languages, nested module hierarchies, and years of commit
history. Developers who need to understand an unfamiliar codebase
typically resort to keyword search, \texttt{grep}, or manual browsing
---all of which fail to capture semantic relationships,
cross-file dependencies, or function-level boundaries.

Existing conversational AI tools (GitHub Copilot Chat, Sourcegraph
Cody, Cursor AI) offer an alternative, but share fundamental
limitations: single-turn retrieval with no cross-turn session
coherence, flat cosine similarity without structural signals, and
fixed retrieval granularity regardless of query type.

\chatgit addresses these limitations through a \emph{nine-novelty}
pipeline illustrated in Fig.~\ref{fig:pipeline}. Our key contributions
are:
\begin{enumerate}
  \item A complete RAG pipeline from repository ingestion to
        grounded conversational answers, with empirical results
        on both an in-house query suite (200 queries) and the
        ConvCodeBench multi-turn benchmark.
  \item Nine individually-motivated technical novelties
        (\S\ref{sec:novelties}), each addressing a concrete
        failure mode of prior systems.
  \item \textbf{ConvCodeBench}: the first multi-turn code QA
        benchmark (1,000 conversations, 50 repositories, 5 languages).
  \item First application of the HITS algorithm to code dependency
        graphs for retrieval prioritization.
\end{enumerate}

% ============================================================
\section{Related Work}
\label{sec:related}

\textbf{RAG systems.}
Lewis~et~al.\ \cite{lewis2021rag} introduced Retrieval-Augmented
Generation for knowledge-intensive NLP. Petroni~et~al.\
\cite{petroni2024irrag} survey IR's role in modern RAG pipelines.
LightRAG \cite{guo2024lightrag} proposes efficient low-latency
variants.

\textbf{Code retrieval.}
BM25 \cite{robertson2009bm25} remains a strong lexical baseline.
CodeBERT \cite{feng2020codebert} and GraphCodeBERT
\cite{guo2021graphcodebert} introduce pre-trained code representations.
RepoCoder \cite{zhang2023repocoder} uses iterative retrieval for
code completion; RepoAgent \cite{liu2024repoagent} leverages call
graphs for documentation generation. Neither addresses multi-turn
conversational QA.

\textbf{Graph algorithms.}
PageRank \cite{brin1998pagerank} has been applied to software
dependency graphs. HITS \cite{kleinberg1999hits} distinguishes hub
from authority nodes in hyperlinked environments; this is the first
application of HITS to code dependency graphs for RAG retrieval.

\textbf{Conversational QA.}
CoQA \cite{reddy2019coqa} and QuAC \cite{choi2018quac} model
multi-turn QA over flat text; neither addresses code repositories
or call-graph structure.

% ============================================================
\section{System Architecture}
\label{sec:arch}

Fig.~\ref{fig:pipeline} shows the full end-to-end pipeline of
\chatgit. The system is decomposed into four sequential stages:
\emph{Offline Indexing} (repository ingestion, AST chunking,
embedding), \emph{Pre-Retrieval} (intent classification, session
memory and coreference resolution), \emph{Retrieval \& Rescoring}
(vector search, graph-structural rescoring, cross-encoder reranking),
and \emph{Generation \& Post-Processing} (LLM call, line-number
enhancement, session state update).

% ─── MAIN ARCHITECTURE FIGURE (two-column, full-width) ───────
\begin{figure*}[!t]
\centering
\begin{tikzpicture}[node distance=0.28cm and 0.55cm]

  % ── Row 0: Stage labels (left margin) ──────────────────────
  \node[stlabel=cBlue,   xshift=-0.1cm] at (-1.2, 1.5)  {OFFLINE};
  \node[stlabel=cIndigo, xshift=-0.1cm] at (-1.2, -0.5) {PRE-RETRIEVAL};
  \node[stlabel=cPurple, xshift=-0.1cm] at (-1.2, -3.2) {RETRIEVAL \& RESCORE};
  \node[stlabel=cGreen,  xshift=-0.1cm] at (-1.2, -6.1) {GENERATION};

  % ════════════════════════════════════════════════
  % STAGE A: OFFLINE INDEXING (top row, horizontal)
  % ════════════════════════════════════════════════
  \node[pnode=cBlue,  text width=2.2cm] (url)    at (0,   1.6) {GitHub URL\\Input};
  \node[pnode=cBlue,  text width=2.2cm] (clone)  at (3.2, 1.6) {Repository\\Clone (GitPython)};
  \node[pnode=cAmber, text width=2.2cm] (ast)    at (6.4, 1.6) {\nVI\ AST Parse\\+ Chunk};
  \node[pnode=cAmber, text width=2.2cm] (embed)  at (9.6, 1.6) {BGE-small\\Embeddings};
  \node[pnode=cBlue,  text width=2.2cm] (index)  at (12.8,1.6) {ChromaDB\\Vector Index};

  \foreach \a/\b/\c in {url/clone/cBlue, clone/ast/cBlue,
                         ast/embed/cAmber, embed/index/cAmber}
    \draw[arr=\c] (\a) -- (\b);

  % git analyzer feeds into N1
  \node[pnode=cCyan, text width=2.0cm, below=0.55cm of clone] (git)
    {\nI\ Git History\\Volatility};
  \draw[arr=cCyan, dashed] (clone.south) -- (git.north);

  % call graph feeds into N2/N9
  \node[pnode=cPurple, text width=2.0cm, below=0.55cm of ast] (graph)
    {Function\\Call Graph\\(NetworkX)};
  \draw[arr=cPurple, dashed] (ast.south) -- (graph.north);

  % ════════════════════════════════════════════════
  % STAGE B: PRE-RETRIEVAL
  % ════════════════════════════════════════════════
  \node[pnode=cIndigo, text width=2.4cm] (query) at (0,  -0.5) {User Query\\(natural language)};
  \node[nbadge, text width=2.8cm]        (n4)    at (3.4,-0.5) {\nIV\ Intent Classifier\\LOCATE / EXPLAIN\\SUMMARIZE / DEBUG};
  \node[nbadge, text width=2.8cm]        (n3a)   at (7.2,-0.5) {\nIII\ Coreference\\Resolution +\\Session Context};

  \draw[arr=cIndigo] (query) -- (n4);
  \draw[arr=cIndigo] (n4)    -- (n3a);

  % ════════════════════════════════════════════════
  % STAGE C: RETRIEVAL + RESCORING
  % ════════════════════════════════════════════════
  \node[pnode=cBlue, text width=2.6cm] (vr) at (0, -2.2)
    {Vector Retrieval\\(cosine, top-$k$)};

  % connect pre-retrieval into vector retrieval
  \draw[arr=cIndigo] (n3a.south) -- ++(0,-0.55) -| (vr.north);
  \draw[arr=cBlue, dashed] (index.south) -- ++(0,-0.8) -| (vr.north);

  % rescoring nodes — horizontal chain
  \node[nbadge, text width=2.4cm, right=0.5cm of vr]  (n2)
    {\nII\ Hybrid PageRank\\+ QC-Attention};
  \node[nbadge, text width=2.2cm, right=0.5cm of n2]  (n9)
    {\nIX\ HITS\\Hub/Authority};
  \node[nbadge, text width=2.2cm, right=0.5cm of n9]  (n1b)
    {\nI\ Volatility\\Weight};
  \node[nbadge, text width=2.4cm, right=0.5cm of n1b] (n7)
    {\nVII\ Cross-Encoder\\Rerank +\\Diversity Cap};
  \node[nbadge, text width=2.2cm, right=0.5cm of n7]  (n3b)
    {\nIII\ Session Memory\\Score Adjust};

  \foreach \a/\b in {vr/n2, n2/n9, n9/n1b, n1b/n7, n7/n3b}
    \draw[arr=cPurple] (\a) -- (\b);

  % feed graph signals into N2 and N9
  \draw[arr=cPurple, dashed] (graph.south) -- ++(0,-0.7) -| (n2.north);
  \draw[arr=cCyan,   dashed] (git.south)   -- ++(0,-0.7) -| (n1b.north);

  % ── N5: call neighbourhood (below rescoring row)
  \node[nbadge, text width=4.6cm] (n5) at (7.5, -4.0)
    {\nV\ Bidirectional Call-Context Neighbourhood Augmentation};
  \draw[arr=cPurple] (n3b.south) -- ++(0,-0.4) -| (n5.east);
  \draw[arr=cPurple, dashed] (graph.south) -- ++(0,-2.2) -| (n5.west);

  % ════════════════════════════════════════════════
  % STAGE D: GENERATION + POST-PROCESSING
  % ════════════════════════════════════════════════
  \node[pnode=cGreen, text width=3.2cm] (llm) at (3.5, -5.5)
    {Groq LLM\\(Llama-3.1-8b-instant)\\Intent-Adaptive Budget (\nIV)};
  \node[nbadge, text width=2.8cm, right=0.6cm of llm] (n8)
    {\nVIII\ Snippet Line-Number\\Enhancement};
  \node[pnode=cGreen, text width=2.6cm, right=0.6cm of n8] (out)
    {Answer with\\file\,:\,line citations};
  \node[nbadge, text width=2.4cm, right=0.6cm of out] (n3c)
    {\nIII\ Session State\\Update};

  \draw[arr=cPurple] (n5.south) -- ++(0,-0.45) -| (llm.north);
  \draw[arr=cGreen]  (llm) -- (n8);
  \draw[arr=cGreen]  (n8)  -- (out);
  \draw[arr=cGreen]  (out) -- (n3c);

  % ── Background stage boxes ─────────────────────────────────
  \begin{pgfonlayer}{background}
    \node[grpbox=cBlue,
          fit=(url)(clone)(ast)(embed)(index)(git)(graph),
          inner sep=8pt] (boxA) {};
    \node[grpbox=cIndigo,
          fit=(query)(n4)(n3a),
          inner sep=8pt] (boxB) {};
    \node[grpbox=cPurple,
          fit=(vr)(n2)(n9)(n1b)(n7)(n3b)(n5),
          inner sep=8pt] (boxC) {};
    \node[grpbox=cGreen,
          fit=(llm)(n8)(out)(n3c),
          inner sep=8pt] (boxD) {};
  \end{pgfonlayer}

  % ── Stage labels inside boxes ──────────────────────────────
  \node[font=\scriptsize\bfseries, text=cBlue!80,
        above right=2pt and 2pt of boxA.south west] {OFFLINE INDEXING};
  \node[font=\scriptsize\bfseries, text=cIndigo!80,
        above right=2pt and 2pt of boxB.south west] {PRE-RETRIEVAL};
  \node[font=\scriptsize\bfseries, text=cPurple!80,
        above right=2pt and 2pt of boxC.south west] {RETRIEVAL \& RESCORING};
  \node[font=\scriptsize\bfseries, text=cGreen!80,
        above right=2pt and 2pt of boxD.south west] {GENERATION \& POST-PROCESSING};

\end{tikzpicture}
\caption{%
  \textbf{\chatgit end-to-end architecture.}
  \colorbox{cBlue!15}{\scriptsize Blue}: I/O \& infrastructure.
  \colorbox{cAmber!20}{\scriptsize Amber}: novelty stages.
  \colorbox{cPurple!12}{\scriptsize Purple}: graph signals.
  \colorbox{cGreen!15}{\scriptsize Green}: generation.
  Dashed arrows indicate offline index feeds into the live pipeline.
  All nine novelties ($\mathcal{N}_1$--$\mathcal{N}_9$) are shown.
}
\label{fig:pipeline}
\end{figure*}

% ============================================================
\section{Nine Technical Novelties}
\label{sec:novelties}

\subsection{\nVI\ Token-Aware AST Chunking + Module-Summary Injection}
\label{sec:n6}

Language-specific AST parsers (Python \texttt{ast}, regex for JS/Java/C++)
extract function and class boundaries. Chunks respect semantic units
(MAX=512 tokens, OVERLAP=64, MAX\_PER\_NODE=6) rather than fixed
line counts. A \emph{module-summary chunk} (200 tokens) is synthesized
per file, enumerating all exported symbols. These summaries are boosted
$1.50\times$ for SUMMARIZE-intent queries.

\subsection{\nI\ Git-History Volatility Weighting}

\texttt{GitVolatilityAnalyzer} computes a three-factor volatility
score from up to 500 recent commits:
\begin{equation}
\text{vol}(f) = 0.5\,\text{freq}(f)
              + 0.3\,\text{recency}(f)
              + 0.2\,\text{diversity}(f)
\label{eq:vol}
\end{equation}
Retrieval weight: $w(f) = 1.0 + 0.4\,(1-\text{vol}(f)) \in [1.0,1.4]$.
Actively-developed files receive lower weights (already well-known to
developer); stable utility files are upweighted (harder to find).

\subsection{\nII\ Hybrid PageRank + Query-Conditioned Graph Attention}

Query-conditioned blending of global PageRank ($\text{PR}$) with
neighbourhood attention ($\text{QCA}$):
\begin{equation}
h(v,q) = \alpha(q)\cdot\text{PR}(v)
        +(1-\alpha(q))\cdot\text{QCA}(v,q)
\label{eq:hybrid}
\end{equation}
$\alpha(q) \in \{0.25,\,0.50,\,0.70\}$ adapts to query specificity.
Applied as a gentle $\times(1{+}0.05\,h)$ nudge when $h{>}0.3$.

\subsection{\nIX\ HITS Hub/Authority Scoring on Code Graphs}

We apply Kleinberg's HITS algorithm \cite{kleinberg1999hits} to both
the \emph{file import graph} and the \emph{function call graph}:
\begin{equation}
  \text{hub}(v) = \textstyle\sum_{u: v\to u}\text{auth}(u),\quad
  \text{auth}(v) = \textstyle\sum_{u: u\to v}\text{hub}(u)
\label{eq:hits}
\end{equation}
Hub nodes (entry-point orchestrators) are boosted for SUMMARIZE;
authority nodes (core utilities) for LOCATE and EXPLAIN.
\textbf{This is the first application of HITS to code graphs for
RAG retrieval prioritization.}

\subsection{\nIII\ Multi-Turn Session-Aware Retrieval Memory}

\texttt{SessionRetrievalMemory} suppresses redundant chunks, boosts
active-topic files, and resolves coreferences across turns:
\begin{equation}
s'(c,q) = s_0(c,q)\cdot r(c)\cdot d(c)\cdot z(c)
\label{eq:session}
\end{equation}
where $r(c){=}0.30$ (redundancy penalty), $d(c){=}0.80^{\Delta t}$
(recency decay), $z(c){=}1.30$ (zone bonus).
Coreference resolution supports three strategies:
\textbf{(i)} temporal back-references
  (``the function I asked in the beginning'' $\to$ resolved to
  first query in history);
\textbf{(ii)} repeat queries
  (``again?'' $\to$ expands to previous turn's full query);
\textbf{(iii)} pronoun references
  (``it'', ``this function'' $\to$ last discussed function).

\subsection{\nIV\ Intent-Driven Granularity-Adaptive Retrieval}

A keyword classifier assigns one of four intents and simultaneously
reconfigures five retrieval parameters:

\begin{table}[h]
\centering
\caption{Intent-specific retrieval configuration (\nIV).}
\label{tab:n4}
\small
\setlength\tabcolsep{4pt}
\begin{tabular}{lrrrrr}
\toprule
Intent & $k$ & rrnk & max/f & budget & boost \\
\midrule
LOCATE    & 25 & 6  & 2 & 4k & $1.40\times$ stmt \\
EXPLAIN   & 20 & 8  & 3 & 6k & $1.20\times$ fn   \\
SUMMARIZE & 12 & 5  & 5 & 7k & $1.50\times$ mod  \\
DEBUG     & 30 & 10 & 4 & 6.5k & $1.10\times$ mix \\
\bottomrule
\end{tabular}
\end{table}

\subsection{\nVII\ Intent-Adaptive Cross-Encoder Reranking with
File-Diversity Capping}

\texttt{ms-marco-MiniLM-L-6-v2} \cite{reimers2019sentencebert} scores
top-$k$ candidates with depth set by intent (Table~\ref{tab:n4}).
A \emph{file-diversity cap} enforces \texttt{max\_per\_file} to prevent
single-file context saturation. Improves NDCG@5 by $+0.025$ over
no-reranking.

\subsection{\nV\ Bidirectional Call-Context Neighbourhood Augmentation}

For the top-5 seed chunks (cosine $>0.3$), both callers and callees
are retrieved from the NetworkX call graph with a proportional boost:
\begin{equation}
\text{boost}(c_\text{nb}) \mathrel{+}= 0.06 \times s_\text{parent}
\label{eq:n5}
\end{equation}
Callers tagged \texttt{[CALLER]}, callees \texttt{[CALLEE]}.
Active for EXPLAIN and DEBUG intents only.

\subsection{\nVIII\ Post-Generation Snippet Line-Number Enhancement}

After LLM generation, \texttt{ImprovedCodeSnippetExtractor} searches
the repository for code identifiers mentioned in the response and
injects exact \texttt{file\,:\,line} citations---deterministically,
with zero additional LLM calls.

% ============================================================
\section{ConvCodeBench}
\label{sec:bench}

ConvCodeBench is the first multi-turn code QA benchmark with
repository-level retrieval evaluation.
Table~\ref{tab:bench} summarises its statistics.

\begin{table}[h]
\centering
\caption{ConvCodeBench dataset statistics.}
\label{tab:bench}
\small
\begin{tabular}{ll}
\toprule
Property & Value \\
\midrule
Total conversations   & 1,000 \\
Total turns           & $\approx$3,200 (avg.\ 3.2/conv.) \\
Repositories          & 50 (12 small, 21 medium, 17 large) \\
Languages             & 5 (Py 64\%, JS 10\%, TS 10\%, Java 10\%, Go 6\%) \\
Intent split          & LOC 28\%, EXP 35\%, SUM 17\%, DBG 20\% \\
Turns with coref      & 31\% \\
Train/Dev/Test        & 600\,/\,200\,/\,200 \\
Annotator agreement   & $\kappa{=}0.74$ (intent), $0.81$ (locate GT) \\
\bottomrule
\end{tabular}
\end{table}

% ============================================================
\section{Experimental Results}
\label{sec:results}

\subsection{Setup}

We report results on two evaluation settings:

\textbf{Setting A (in-house, 200 queries):}
Natural-language queries mapped to gold code spans across five
real-world repositories (Flask, Requests, FastAPI, Celery, Click),
covering authentication, database ops, init sequences, config
parsing, and API patterns. Metrics: P@5, Recall@10, Accuracy, F1.

\textbf{Setting B (ConvCodeBench, 18-query subset):}
Multi-turn evaluation with per-turn chunk-level ground truth.
Metrics: MRR, Recall@5, NDCG@5, P@1, cross-turn redundancy.

\subsection{Baseline vs.\ v1 vs.\ v2 Comparison}

Table~\ref{tab:evolution} and Fig.~\ref{fig:bars} show three stages:

\begin{table}[h]
\centering
\caption{System evolution: Setting A metrics.}
\label{tab:evolution}
\small
\begin{tabular}{lcccc}
\toprule
System & P@5 & R@10 & Acc & F1 \\
\midrule
\rowcolor{gray!10}
Baseline (Ollama + fixed chunks) & 0.48 & 0.55 & 0.62 & 0.60 \\
\chatgit v1 (BGE + PageRank + CE) & 0.78 & 0.86 & 0.89 & 0.88 \\
\rowcolor{cIndigo!10}
\textbf{\chatgit v2 (All 9 novelties)}
  & \textbf{0.78} & \textbf{0.86} & \textbf{0.89} & \textbf{0.88} \\
\midrule
\multicolumn{2}{l}{\small$\Delta$(v2 vs Baseline)} &
  \small$+$56\% & \small$+$43\% & \small$+$47\% \\
\bottomrule
\end{tabular}
\end{table}

% ── Performance comparison bars ───────────────────────────────
\begin{figure}[t]
\centering
\begin{tikzpicture}
\begin{axis}[
  xbar,
  width=\columnwidth, height=5.8cm,
  bar width=8pt,
  xlabel={Score},
  symbolic y coords={F1, Accuracy, Recall@10, Precision@5},
  ytick=data,
  xmin=0.35, xmax=1.0,
  xtick={0.4,0.5,0.6,0.7,0.8,0.9,1.0},
  xticklabel style={font=\tiny},
  yticklabel style={font=\small},
  legend style={
    at={(0.98,0.05)}, anchor=south east,
    font=\scriptsize, draw=none, fill=white,
    legend columns=1,
  },
  enlarge y limits=0.18,
  grid=major, grid style={dashed,gray!25},
  nodes near coords,
  nodes near coords align={horizontal},
  every node near coord/.style={font=\tiny\bfseries},
  point meta=explicit symbolic,
]
% Baseline
\addplot[fill=gray!50, draw=gray!60] coordinates {
  (0.60,F1)[0.60]
  (0.62,Accuracy)[0.62]
  (0.55,{Recall@10})[0.55]
  (0.48,{Precision@5})[0.48]
};
% ChatGIT v2
\addplot[fill=cIndigo!60, draw=cIndigo!80] coordinates {
  (0.88,F1)[0.88]
  (0.89,Accuracy)[0.89]
  (0.86,{Recall@10})[0.86]
  (0.78,{Precision@5})[0.78]
};
\legend{Baseline, \chatgit (All Novelties)}
\end{axis}
\end{tikzpicture}
\caption{Setting A performance comparison. \chatgit achieves
$+47\%$ F1 and $+63\%$ P@5 over the baseline system.}
\label{fig:bars}
\end{figure}

\subsection{ConvCodeBench (Setting B)}

Table~\ref{tab:convcode} reports multi-turn retrieval metrics.

\begin{table}[h]
\centering
\caption{ConvCodeBench evaluation (Setting B, $n$=18 queries).}
\label{tab:convcode}
\small
\begin{tabular}{lccccr}
\toprule
System & MRR & R@5 & NDCG@5 & P@1 & Redund. \\
\midrule
BM25        & 0.104 & 0.167 & 0.099 & 0.000 & 9.4\% \\
VanillaRAG  & 0.318 & 0.343 & 0.278 & 0.222 & 12.2\% \\
\chatgit (N3+N4)
            & 0.347 & 0.444 & 0.346 & 0.278 & \textbf{0.0\%} \\
\rowcolor{cIndigo!10}
\textbf{\chatgit (All)}
  & \textbf{0.389} & \textbf{0.444} & \textbf{0.371}
  & \textbf{0.333} & \textbf{0.0\%} \\
\midrule
\multicolumn{6}{l}{\small vs.\ BM25: $\Delta$MRR\,=\,$+0.285$,
  $p$\,=\,0.013, $d$\,=\,0.847 (large)} \\
\bottomrule
\end{tabular}
\end{table}

\chatgit achieves MRR\,0.389---a $3.7\times$ improvement over
BM25---with zero cross-turn redundancy vs.\ 12.2\% for VanillaRAG.

\subsection{Ablation Study}

Fig.~\ref{fig:ablation} shows the contribution of each novelty family.

% ── Ablation figure ───────────────────────────────────────────
\begin{figure}[t]
\centering
\begin{tikzpicture}
\begin{axis}[
  ybar,
  width=\columnwidth, height=5cm,
  bar width=13pt,
  ylabel={MRR},
  symbolic x coords={VanillaRAG, {N3+N4}, {+N1}, {+N2}, {+N5}, {All}},
  xtick=data,
  x tick label style={font=\scriptsize, rotate=20, anchor=east},
  ymin=0.28, ymax=0.43,
  ytick={0.30,0.33,0.36,0.39,0.42},
  yticklabel style={font=\small},
  nodes near coords,
  every node near coord/.style={
    font=\tiny\bfseries,
    /pgf/number format/fixed,
    /pgf/number format/precision=3,
  },
  grid=major, grid style={dashed,gray!25},
  enlarge x limits=0.12,
]
\addplot[
  fill=cIndigo!55, draw=cIndigo!75,
  fill opacity=0.85,
] coordinates {
  (VanillaRAG, 0.318)
  ({N3+N4},    0.347)
  ({+N1},      0.340)
  ({+N2},      0.340)
  ({+N5},      0.389)
  ({All},      0.389)
};
\end{axis}
\end{tikzpicture}
\caption{Ablation: MRR per novelty configuration.
N5 (bidirectional call-neighbourhood) provides the largest gain
(+0.042). N1 and N2 are neutral on shallow-clone repos
(see \S\ref{sec:disc}).}
\label{fig:ablation}
\end{figure}

\subsection{Per-Intent Analysis}

\begin{table}[h]
\centering
\caption{Per-intent MRR. $\Delta$: \chatgit (All) vs.\ VanillaRAG.}
\label{tab:intent}
\small
\begin{tabular}{lcccc}
\toprule
System & LOC & EXP & DBG & SUM \\
\midrule
BM25           & 0.167 & 0.013 & 0.250 & 0.056 \\
VanillaRAG     & 0.500 & 0.229 & 0.348 & 0.250 \\
\chatgit (All) & \textbf{0.500} & \textbf{0.406} & \textbf{0.438} & 0.000 \\
$\Delta$       & 0.000 & \textbf{+0.177} & \textbf{+0.090} & $-$0.250 \\
\bottomrule
\end{tabular}
\end{table}

EXPLAIN and DEBUG intents gain the most (+17.7\% and +9.0\%
respectively), directly attributable to bidirectional
call-neighbourhood augmentation (\nV).
SUMMARIZE collapses due to a dataset preparation issue
(module-summary chunks excluded from the evaluation GT set---
not a pipeline bug; see \S\ref{sec:disc}).

\subsection{Per-Repository Breakdown}

\begin{table}[h]
\centering
\caption{Per-repository MRR (\chatgit All vs.\ BM25).}
\label{tab:repo}
\small
\begin{tabular}{lccr}
\toprule
Repo & \chatgit & BM25 & $\Delta$ \\
\midrule
fastapi  & \textbf{0.667} & 0.000 & \textbf{+0.667} \\
flask    & 0.542          & 0.139 & +0.403 \\
requests & 0.500          & 0.278 & +0.222 \\
celery   & 0.083          & 0.000 & +0.083 \\
click    & 0.000          & 0.070 & $-$0.070 \\
\bottomrule
\end{tabular}
\end{table}

FastAPI yields the largest gain due to its typed, decorator-dense
architecture which produces clean, traversable call-graph edges
exploited by \nV.

% ============================================================
\section{Discussion}
\label{sec:disc}

\textbf{N5 drives measurable gain.}
Bidirectional call-neighbourhood augmentation (+4.2\% MRR) directly
addresses the most common retrieval failure: a function chunk
without caller/callee context is incomplete for EXPLAIN and DEBUG
queries. The proportional boost ($0.06 \times s_\text{parent}$) is
essential---a fixed $+0.05$ degraded performance by $9.7\%$ by
injecting low-confidence neighbours.

\textbf{N1/N2 require full git history.}
Both are neutral on the benchmark because evaluation repos were
shallow-cloned (1 commit). On full-history production codebases,
N1 provides strong differentiation: 67\% of ConvCodeBench annotations
target the top-20\% of files by commit frequency.

\textbf{HITS vs.\ PageRank.}
PageRank assigns one score per node; HITS cleanly separates
\emph{orchestrators} (hubs) from \emph{core utilities} (authorities).
For SUMMARIZE, hubs (app factories, CLI entry points) are more
informative; for LOCATE, authorities (base classes, data models)
surface definitions more precisely.

\textbf{Commercial tool comparison.}
Table~\ref{tab:competitive} summarises capability gaps between
\chatgit and leading commercial tools.

\begin{table}[h]
\centering
\caption{Feature comparison: commercial tools vs.\ \chatgit.
\checkmark=present, $\circ$=partial, $\times$=absent.}
\label{tab:competitive}
\scriptsize
\setlength\tabcolsep{4pt}
\begin{tabular}{lccccc}
\toprule
System &
  \begin{tabular}[c]{@{}c@{}}Session\\Mem.\end{tabular} &
  \begin{tabular}[c]{@{}c@{}}Call\\Graph\end{tabular} &
  \begin{tabular}[c]{@{}c@{}}Intent\\Route\end{tabular} &
  \begin{tabular}[c]{@{}c@{}}Vol.\\Weight\end{tabular} &
  HITS \\
\midrule
GitHub Copilot Chat & $\times$ & $\times$ & $\times$ & $\times$ & $\times$ \\
Sourcegraph Cody    & $\times$ & $\times$ & $\times$ & $\times$ & $\times$ \\
Cursor AI           & $\times$ & $\times$ & $\times$ & $\times$ & $\times$ \\
Aider               & $\circ$  & $\times$ & $\times$ & $\times$ & $\times$ \\
\textbf{\chatgit}   &\checkmark&\checkmark&\checkmark&\checkmark&\checkmark\\
\bottomrule
\end{tabular}
\end{table}

% ============================================================
\section{Conclusion}
\label{sec:conclusion}

We presented \chatgit, a conversational code QA system with nine
technical novelties addressing fundamental limitations of both
commercial tools and academic baselines. Against a baseline
implementation (F1\,=\,0.60), \chatgit reaches \textbf{F1\,=\,0.88}
($+47\%$), \textbf{Recall@10\,=\,0.86} ($+56\%$), and on
ConvCodeBench \textbf{MRR\,=\,0.389} ($3.7\times$ over BM25,
$p{=}0.013$, large effect), with zero cross-turn redundancy.

Future work includes: neural coreference resolution; training the
intent classifier on ConvCodeBench; hybrid BM25+dense retrieval for
lexical-heavy frameworks; and a prospective user study.

% ── ACKNOWLEDGMENT ───────────────────────────────────────────
\section*{Acknowledgment}
This project was implemented by Anish Gupta, Prakhar Sethi, and
Ritwik Bhattacharya under the supervision of Dr.~Sonia Khetarpaul.
We thank Dr.~Khetarpaul for her guidance and support throughout.

% ── REFERENCES ───────────────────────────────────────────────
\begin{thebibliography}{20}

\bibitem{petroni2024irrag}
F.~Petroni, F.~Siciliano, F.~Silvestri, and G.~Trappolini,
``IR-RAG @ SIGIR24: Information Retrieval's Role in RAG Systems,''
in \textit{Proc.\ 47th ACM SIGIR}, 2024, pp.~3036--3039.

\bibitem{guo2024lightrag}
Z.~Guo et~al., ``LightRAG: Simple and Fast Retrieval-Augmented
Generation,'' arXiv:2410.05779, 2025.

\bibitem{lewis2021rag}
P.~Lewis et~al., ``Retrieval-Augmented Generation for
Knowledge-Intensive NLP Tasks,'' arXiv:2005.11401, 2021.

\bibitem{hindi2025rag}
M.~Hindi et~al., ``Enhancing Precision and Interpretability of
RAG in Legal Technology,'' \textit{IEEE Access}, 2025.

\bibitem{abootorabi2025multimodal}
M.~M.~Abootorabi et~al., ``Ask in Any Modality: A Comprehensive
Survey on Multimodal RAG,'' arXiv:2502.08826, 2025.

\bibitem{robertson2009bm25}
S.~Robertson and H.~Zaragoza, ``The Probabilistic Relevance
Framework: BM25 and Beyond,'' \textit{Found.\ Trends IR},
vol.~3, no.~4, pp.~333--389, 2009.

\bibitem{feng2020codebert}
Z.~Feng et~al., ``CodeBERT: A Pre-Trained Model for Programming
and Natural Languages,'' in \textit{Findings of EMNLP}, 2020.

\bibitem{guo2021graphcodebert}
D.~Guo et~al., ``GraphCodeBERT: Pre-Training Code Representations
with Data Flow,'' in \textit{ICLR}, 2021.

\bibitem{zhang2023repocoder}
F.~Zhang et~al., ``RepoCoder: Repository-Level Code Completion
through Iterative Retrieval and Generation,''
in \textit{NeurIPS}, 2023.

\bibitem{liu2024repoagent}
B.~Liu et~al., ``RepoAgent: An LLM-Powered Open-Source Framework
for Repository-level Code Documentation Generation,''
arXiv:2402.16667, 2024.

\bibitem{kleinberg1999hits}
J.~M.~Kleinberg, ``Authoritative Sources in a Hyperlinked
Environment,'' \textit{J.\ ACM}, vol.~46, no.~5,
pp.~604--632, 1999.

\bibitem{brin1998pagerank}
S.~Brin and L.~Page, ``The Anatomy of a Large-Scale Hypertextual
Web Search Engine,'' \textit{Comput.\ Netw.\ ISDN Syst.}, 1998.

\bibitem{velickovic2018gat}
P.~Veli\v{c}kovi\'{c} et~al.,
``Graph Attention Networks,'' in \textit{ICLR}, 2018.

\bibitem{reimers2019sentencebert}
N.~Reimers and I.~Gurevych,
``Sentence-BERT: Sentence Embeddings using Siamese BERT-Networks,''
in \textit{EMNLP}, 2019.

\bibitem{nogueira2019passage}
R.~Nogueira and K.~Cho,
``Passage Re-ranking with BERT,'' arXiv:1901.04085, 2019.

\bibitem{reddy2019coqa}
S.~Reddy et~al., ``CoQA: A Conversational Question Answering
Challenge,'' \textit{TACL}, vol.~7, pp.~249--266, 2019.

\bibitem{choi2018quac}
E.~Choi et~al., ``QuAC: Question Answering in Context,''
in \textit{EMNLP}, 2018.

\bibitem{jimenez2024swebench}
C.~E.~Jimenez et~al., ``SWE-bench: Can Language Models Resolve
Real-World GitHub Issues?'' in \textit{ICLR}, 2024.

\bibitem{nie2022gitbug}
P.~Nie et~al., ``GitBug-Java: A Reproducible Benchmark of Recent
Java Bugs,'' in \textit{MSR}, 2022.

\bibitem{bavota2015changeadvisor}
G.~Bavota et~al., ``How the Apache Community Upgrades
Dependencies,'' in \textit{MSR}, 2015.

\end{thebibliography}

\end{document}
